% Осень26!A26; нумерация внутри столбца без пустых строк.
% Источник: лекция 4, §1 (строка 27)
\ticket{23}{Пример неизмеримого по Лебегу множества}
\label{ticket:23}

\begin{theorem}
Существует множество $A\subset[0,1]$, не измеримое по Лебегу.
\end{theorem}

\begin{proof}
На $[0,1]$ введём отношение $x\sim y\iff x-y\in\mathbb Q$.
По аксиоме выбора возьмём из каждого класса по одному представителю и обозначим их множество через $A$.
Для $Q=\mathbb Q\cap[-1,1]$ имеем
\[
  [0,1]\subset\bigcup_{q\in Q}(A+q)\subset[-1,2].
\]
Действительно, для каждого $x\in[0,1]$ найдётся представитель $a\in A$ того же класса, и $q=x-a\in Q$.
Сдвиги $A+q$ попарно не пересекаются: из $a_1+q_1=a_2+q_2$ следует $a_1\sim a_2$, затем $a_1=a_2$ и $q_1=q_2$.

Предположим, что $A$ измеримо.
Используем инвариантность измеримости и меры Лебега относительно сдвигов.
По счётной аддитивности
\[
  1\le\lambda\left(\bigsqcup_{q\in Q}(A+q)\right)
    =\sum_{q\in Q}\lambda(A)\le3.
\]
Если $\lambda(A)=0$, сумма равна нулю; если $\lambda(A)>0$, она бесконечна.
Оба случая противоречат полученным оценкам.
\end{proof}
